\begin{center}
    \textbf{Vilniaus Universitetas} \\
    \textbf{Sistemų Biologija} \\
    
    \vspace{0.5cm}
    
    \textbf{PAVIENIŲ LĄSTELIŲ RNR SEKOSKAITOS DUOMENŲ ANALIZĖS TOBULINIMAS INTEGRUOJANT NEAPRAŠYTUS TARPGENINIUS REGIONUS} \\
    
    \vspace{1cm}
    
    \textbf{Juozapas Ivanauskas} \\
    
    \vspace{1cm}
    
    \textbf{Santrauka} \\
    
\end{center}

Pavienių ląstelių RNR sekoskaitos eksperimentai sugeneruoja didelius duomenų kiekius, kurių dalį sudaro sekos, prilygiuojamos prie genomo,
bet nepriskiriaamos jokiam genui.
Šio projekto tikslas buvo ištirti šias sekas ir pabandyti įtraukti jas į tolimesnes analizes.

Šiame projekte buvo išanalizuota nepriskirtos sekos iš keturiolikos duomenų rinkinių.
Dauguma jų buvo priligiuojamos prie priešingos DNR grandies nei žinomi genai.
Buvo iškelta hipotezė, kad šios sekos yra artefaktai susidarantys atvirkštinės transkripcijos metu
(ši hipotezė galioja tik sekoskaitos metodams naudojantiems poli-T pradmenis).
Kita dalis sekų buvo priligiuota prie genomo vietose, kuriose iš viso nėra žinomų genų.
Šios sekos galimai kyla iš neaprašytų genų.
147 regionai buvo iškelti kaip kandidatai naujiems genams, kuriuos reikėtų patvirtinti papildomais ekperimentais.

Galiausiai, dalis sekų buvo priligiuojamos tose vietose, kur yra žinomų genų, tačiau vis tiek liko nepriskirtos.
Kad įtrauktume šias sekas į tolimesnę analizę, patobulinome transkriptominį anotacijos failą.
Tai buvo pasiekta išsprendžiant persidengiančius genus,
pratęsiant kai kurių genų 3' galus bei naudojant informaciją iš kelių skirtingų transkriptomo anotacijų.
Taikant šį metodą, duomenų kiekis tolesnėje analizėje padidėjo vidutiniškai 1.2\%.

Šie rezultatai rodo, kad nepriskiriamos sekos (įprastai neįtraukiamos į tolimesnes analizes), gali būti informatyvios ir
turėtų būti sistemiškai išanalizuotos siekiant pagerinti dabartines genomo anotacijas bei pavienių ląstelių duomenų analizę.
